\documentclass[11pt,a4paper]{article}

\usepackage[T1]{fontenc}
\usepackage[utf8]{inputenc}
\usepackage{lmodern}
\usepackage{amssymb,amsmath,amsthm,mathtools}
\usepackage[margin=2.5cm]{geometry}
\usepackage{parskip}
\usepackage[hidelinks]{hyperref}
\usepackage{microtype}

% Theorem environments (arXiv / journal style)
\theoremstyle{plain}
\newtheorem{theorem}{Theorem}[section]
\newtheorem{proposition}[theorem]{Proposition}
\newtheorem{lemma}[theorem]{Lemma}
\newtheorem{corollary}[theorem]{Corollary}

\theoremstyle{definition}
\newtheorem{definition}[theorem]{Definition}
\newtheorem{example}[theorem]{Example}
\newtheorem{remark}[theorem]{Remark}

% Notation
\newcommand{\Type}{\mathbf{Type}}
\newcommand{\cmin}{\mathrm{cmin}}
\newcommand{\cmax}{\mathrm{cmax}}
\newcommand{\tauu}{\tau}

\title{Trajectory Systems:\\
  Circular Verification Collapses Canonical Trust}
\author{Extracted from a machine-checked development}
\date{}

\begin{document}
\maketitle

\begin{abstract}
A trajectory system is a trust system in which claims are agent actions
and checking steps are causal dependencies among those actions. We introduce
a five-level trajectory chain and two concrete systems sharing a common
dependency topology: a \emph{circular} trajectory, modelling circular
verification, and an \emph{honest} trajectory with sound grades. In the
circular system, canonical trust $\tauu$ evaluates to bottom at every event.
In the honest system, $\tauu$ attains the grade $\mathsf{independent}$ along
sound paths and is limited by the weakest step grade along weaker ones. The
proofs of all stated results have been machine-checked; we record only the
mathematical content.
\end{abstract}

\tableofcontents
\bigskip

%======================================================================
\section{Introduction}
%======================================================================

Trust systems grade claims by the strength of the checks that support them.
Canonical trust $\tauu$ is the greatest sound assignment of grades: at each
claim, the best (join) over checking steps of the meet of the step's own
grade with the canonical trusts of its premises. When the premise relation
is well founded, $\tauu$ is uniquely determined by well-founded recursion.

Trajectory systems specialise this picture to agent workflows. Claims are
events (actions of an agent); checking steps record causal dependence of one
event on others; grades record how strongly each dependence was checked.
The motivating phenomenon is \emph{circular verification}: a chain of events
in which each step claims independence of prior checks, yet the root is
unsupported and a terminal ``verify'' step is itself unsupported. The
widest-bottleneck semantics of $\tauu$ then forces every event to bottom,
regardless of intermediate grades that look strong in isolation.

We work throughout with finite chains of check strengths and finite systems
of events. No special order theory beyond meets and joins in a chain is
required. Proofs are indicated by the phrase ``by lean,'' meaning they have
been checked by a proof assistant rather than written out in prose.

\paragraph{What this note claims.}
In ordinary mathematical language, the claim is: under the circular grading
of Definition~\ref{def:circular}, canonical trust is identically bottom; under
the honest grading of Definition~\ref{def:honest} on the same dependency
topology, canonical trust recovers the independent grade along the sound
spine and is limited by the shallow grade at the weaker branch. Circular
verification is therefore invisible to node-local grade inspection alone: it
is diagnosed only by computing $\tauu$ globally.

%======================================================================
\section{Recollections: chains and trust systems}
%======================================================================

We recall only what is needed; full development appears in the companion
note on trust systems.

\begin{definition}[Chain]\label{def:chain}
A \textbf{chain} $T$ is a linearly ordered type equipped with a bottom
element $\bot$, such that $\le$ is reflexive, transitive, antisymmetric, and
total, and $\bot \le a$ for every $a \in T$. Meets and joins exist and are
computable:
\begin{equation}
  \cmin(a,b) \;=\;
  \begin{cases}
    a & \text{if } a \le b, \\
    b & \text{otherwise,}
  \end{cases}
  \qquad
  \cmax(a,b) \;=\;
  \begin{cases}
    b & \text{if } a \le b, \\
    a & \text{otherwise.}
  \end{cases}
\end{equation}
\end{definition}

\begin{definition}[Trust system]\label{def:system}
A \textbf{trust system} over claims $C$ graded in a chain $T$ consists of
\begin{enumerate}
  \item a type $\mathrm{Step}$ of checking steps,
  \item a conclusion map $\mathrm{concl} : \mathrm{Step} \to C$,
  \item a finite list of premises $\mathrm{prems} : \mathrm{Step} \to
        \mathrm{List}(C)$,
  \item a grade map $\mathrm{grade} : \mathrm{Step} \to T$,
  \item for each claim $x$, a finite list $\mathrm{stepsInto}(x)$ of steps
        concluding at $x$, coherent with $\mathrm{concl}$.
\end{enumerate}
The \textbf{premise relation} $y \prec x$ holds when $y$ is a premise of
some step into $x$.
\end{definition}

\begin{definition}[Step grade and canonical trust]\label{def:tau}
Given an assignment $t : C \to T$, the \textbf{step grade} of a step $f$ is
\begin{equation}
  \mathrm{stepGrade}(t,f)
  \;=\;
  \cmin\bigl(
    \mathrm{grade}(f),\;
    \cmin_{y \in \mathrm{prems}(f)} t(y)
  \bigr)
\end{equation}
(the meet of the step's own grade with the assigned levels of all premises;
empty meets yield the step grade alone).

When the premise relation is well founded, \textbf{canonical trust}
$\tauu : C \to T$ is defined by well-founded recursion:
\begin{equation}
  \tauu(x)
  \;=\;
  \cmax\Bigl(
    \bot,\;
    \cmax_{f \in \mathrm{stepsInto}(x)}
      \mathrm{stepGrade}(\tauu, f)
  \Bigr).
\end{equation}
Equivalently, $\tauu(x)$ is the join over steps into $x$ of their step
grades under $\tauu$ itself, defaulting to $\bot$ when there are no steps.
\end{definition}

\begin{remark}
Canonical trust is the greatest sound assignment: every sound grading is
bounded by $\tauu$, and $\tauu$ itself is sound. We use only the recursive
character of $\tauu$ below.
\end{remark}

%======================================================================
\section{The trajectory chain}
%======================================================================

\begin{definition}[Trajectory levels]\label{def:tlevel}
The \textbf{trajectory chain} $T_{\mathrm{traj}}$ is the five-element chain
\begin{equation}
  \mathsf{unsupported}
  \;<\;
  \mathsf{shallow}
  \;<\;
  \mathsf{independent}
  \;<\;
  \mathsf{crossVerified}
  \;<\;
  \mathsf{mechChecked},
\end{equation}
with bottom $\bot = \mathsf{unsupported}$. The order is the unique linear
order extending the displayed chain of covering relations (and equality).
\end{definition}

\begin{proposition}[Chain structure]\label{prop:tlevel-chain}
$T_{\mathrm{traj}}$ carries the structure of a chain in the sense of
Definition~\ref{def:chain}.
\end{proposition}

\begin{proof}
By lean.
\end{proof}

\begin{remark}
Informally: $\mathsf{unsupported}$ means no usable check;
$\mathsf{shallow}$ a weak or one-sided check;
$\mathsf{independent}$ a check that does not depend on the claim under
review; $\mathsf{crossVerified}$ independent confirmation from a second
source; $\mathsf{mechChecked}$ machine verification.
\end{remark}

%======================================================================
\section{Events and shared topology}
%======================================================================

\begin{definition}[Events]\label{def:evt}
Let $E = \{e_0, e_1, e_2, e_3, e_4\}$ be a five-element type of
\textbf{events}.
\end{definition}

Both systems below use the same dependency topology: steps are identified
with events (one step concluding at each event), and premises form a
directed acyclic graph
\begin{equation}
  e_0 \;\longrightarrow\; e_1 \;\longrightarrow\; e_2
  \;\longrightarrow\; e_3,
  \qquad
  e_2 \;\longrightarrow\; e_4.
\end{equation}
Explicitly:
\begin{center}
\begin{tabular}{lll}
\hline
Event & Premises & Informal role \\
\hline
$e_0$ & $\emptyset$ & root / ground event \\
$e_1$ & $\{e_0\}$ & first dependent step \\
$e_2$ & $\{e_1\}$ & intermediate step \\
$e_3$ & $\{e_2\}$ & claim branch \\
$e_4$ & $\{e_2\}$ & verify branch \\
\hline
\end{tabular}
\end{center}

\begin{proposition}[Well-foundedness of the topology]
The premise relation of this topology is well founded.
\end{proposition}

\begin{proof}
By lean.
\end{proof}

%======================================================================
\section{The circular trajectory}
%======================================================================

\begin{definition}[Circular trajectory]\label{def:circular}
The \textbf{circular trajectory} is the trust system
$\mathcal{S}_{\mathrm{circ}}$ on claims $E$ graded in $T_{\mathrm{traj}}$
with the topology of the previous section and grades
\begin{center}
\begin{tabular}{ll}
\hline
Event & Grade \\
\hline
$e_0$ & $\mathsf{unsupported}$ \\
$e_1$ & $\mathsf{independent}$ \\
$e_2$ & $\mathsf{independent}$ \\
$e_3$ & $\mathsf{shallow}$ \\
$e_4$ & $\mathsf{unsupported}$ \\
\hline
\end{tabular}
\end{center}
\end{definition}

\begin{remark}[Why ``circular'']
The topology is acyclic; circularity lives in the grades. The root $e_0$ is
unsupported, yet $e_1$ and $e_2$ are graded as independent checks, and the
terminal verify step $e_4$ is itself unsupported. Locally, intermediate
steps look strong; globally, every path either begins at an unsupported
root or ends at an unsupported verify. The computation of $\tauu$ makes this
precise.
\end{remark}

\begin{lemma}[Leaf computation]\label{lem:circ-e0}
In $\mathcal{S}_{\mathrm{circ}}$,
\begin{equation}
  \tauu(e_0) \;=\; \mathsf{unsupported}.
\end{equation}
\end{lemma}

\begin{proof}
By lean.
\end{proof}

\begin{lemma}\label{lem:circ-e1}
In $\mathcal{S}_{\mathrm{circ}}$,
\begin{equation}
  \tauu(e_1) \;=\; \mathsf{unsupported}.
\end{equation}
\end{lemma}

\begin{proof}
By lean.
\end{proof}

\begin{lemma}\label{lem:circ-e2}
In $\mathcal{S}_{\mathrm{circ}}$,
\begin{equation}
  \tauu(e_2) \;=\; \mathsf{unsupported}.
\end{equation}
\end{lemma}

\begin{proof}
By lean.
\end{proof}

\begin{lemma}\label{lem:circ-e3}
In $\mathcal{S}_{\mathrm{circ}}$,
\begin{equation}
  \tauu(e_3) \;=\; \mathsf{unsupported}.
\end{equation}
\end{lemma}

\begin{proof}
By lean.
\end{proof}

\begin{lemma}\label{lem:circ-e4}
In $\mathcal{S}_{\mathrm{circ}}$,
\begin{equation}
  \tauu(e_4) \;=\; \mathsf{unsupported}.
\end{equation}
\end{lemma}

\begin{proof}
By lean.
\end{proof}

\begin{theorem}[Circular collapse]\label{thm:circular-bot}
In the circular trajectory, canonical trust is bottom at every event:
\begin{equation}
  \forall\, e \in E,\qquad
  \tauu_{\mathcal{S}_{\mathrm{circ}}}(e) \;=\; \mathsf{unsupported}.
\end{equation}
\end{theorem}

\begin{proof}
By lean.
\end{proof}

\begin{remark}[Idea of the argument]
Proceed leaf-to-root along the premise order. At $e_0$ the unique step has
empty premises and grade $\mathsf{unsupported}$, so
$\tauu(e_0) = \mathsf{unsupported}$. At $e_1$ the step grade is
$\cmin(\tauu(e_0), \mathsf{independent}) = \mathsf{unsupported}$, and
similarly for $e_2$. At $e_3$ one meets $\tauu(e_2)$ with
$\mathsf{shallow}$; at $e_4$ one meets $\tauu(e_2)$ with
$\mathsf{unsupported}$. In every case the join with $\bot$ remains
$\mathsf{unsupported}$. Intermediate grades of $\mathsf{independent}$ never
appear in $\tauu$ because they are bottlenecked by unsupported ancestors
(or, for $e_4$, by an unsupported self-grade).
\end{remark}

%======================================================================
\section{The honest trajectory}
%======================================================================

\begin{definition}[Honest trajectory]\label{def:honest}
The \textbf{honest trajectory} is the trust system
$\mathcal{S}_{\mathrm{hon}}$ on the same claims $E$, the same topology, and
grades
\begin{center}
\begin{tabular}{ll}
\hline
Event & Grade \\
\hline
$e_0$ & $\mathsf{independent}$ \\
$e_1$ & $\mathsf{independent}$ \\
$e_2$ & $\mathsf{independent}$ \\
$e_3$ & $\mathsf{shallow}$ \\
$e_4$ & $\mathsf{independent}$ \\
\hline
\end{tabular}
\end{center}
\end{definition}

\begin{proposition}[Well-foundedness]
The premise relation of $\mathcal{S}_{\mathrm{hon}}$ is well founded.
\end{proposition}

\begin{proof}
By lean.
\end{proof}

\begin{theorem}[Honest root]\label{thm:hon-e0}
In $\mathcal{S}_{\mathrm{hon}}$,
\begin{equation}
  \tauu(e_0) \;=\; \mathsf{independent}.
\end{equation}
\end{theorem}

\begin{proof}
By lean.
\end{proof}

\begin{theorem}[Honest spine]\label{thm:hon-e1}
In $\mathcal{S}_{\mathrm{hon}}$,
\begin{equation}
  \tauu(e_1) \;=\; \mathsf{independent}.
\end{equation}
\end{theorem}

\begin{proof}
By lean.
\end{proof}

\begin{theorem}\label{thm:hon-e2}
In $\mathcal{S}_{\mathrm{hon}}$,
\begin{equation}
  \tauu(e_2) \;=\; \mathsf{independent}.
\end{equation}
\end{theorem}

\begin{proof}
By lean.
\end{proof}

\begin{theorem}[Weak claim branch]\label{thm:hon-e3}
In $\mathcal{S}_{\mathrm{hon}}$,
\begin{equation}
  \tauu(e_3) \;=\; \mathsf{shallow}.
\end{equation}
\end{theorem}

\begin{proof}
By lean.
\end{proof}

\begin{theorem}[Honest verify branch]\label{thm:hon-e4}
In $\mathcal{S}_{\mathrm{hon}}$,
\begin{equation}
  \tauu(e_4) \;=\; \mathsf{independent}.
\end{equation}
\end{theorem}

\begin{proof}
By lean.
\end{proof}

\begin{remark}[Contrast with the circular case]
The topology is identical; only the grades differ. With an independent root
and independent steps along the spine $e_0 \to e_1 \to e_2 \to e_4$,
canonical trust preserves $\mathsf{independent}$. The claim branch $e_3$ is
limited by its own shallow grade: weakest-link semantics yields
$\cmin(\tauu(e_2), \mathsf{shallow}) = \mathsf{shallow}$. Thus $\tauu$
distinguishes sound verification structure from circular structure on the
nose.
\end{remark}

%======================================================================
\section{Discussion}
%======================================================================

\paragraph{What has been established.}
Trajectory systems are trust systems specialised to agent-action claims.
On a fixed five-event topology, circular grades force
$\tauu \equiv \mathsf{unsupported}$, while honest grades recover
$\mathsf{independent}$ along the sound path and $\mathsf{shallow}$ on the
weak branch. The diagnosis of circular verification is therefore a theorem
about $\tauu$, not a side condition on individual step grades.

\paragraph{Role of well-foundedness.}
Both systems have well-founded premise relations (the dependency graph is a
finite DAG). Well-foundedness is load-bearing for the existence of $\tauu$ as
a recursive definition; circular \emph{citation} in the sense of a cyclic
premise relation is a different pathology (treated in the companion note on
trust systems). Here the circularity is graded rather than structural.

\paragraph{Scope.}
The development treats a single finite event set and a single five-level
chain. It does not address infinite trajectories, continuous time, or
composition of trajectory systems with other trust systems; the latter is
taken up in a companion note on trajectory--trust morphisms.

\paragraph{Machine checking.}
Every definition and theorem statement above corresponds to a declaration in
a machine-checked formalisation. The phrase ``by lean'' indicates that the
proof has been verified by a proof assistant; we have not reproduced the
formal proof scripts here. A mathematician reading this note need only
verify that the English statements match the intended mathematics; the
formal system is responsible for the derivations.

%======================================================================
\section*{Acknowledgements}
%======================================================================

The mathematical content specialises the theory of trust systems (canonical
trust as greatest sound assignment, well-founded premise recursion) to
trajectory workflows. The circular example is motivated by observed agent
behaviour under circular verification. The presentation here is a prose
rendering of a formal development intended for human review of the theorem
statements.

\begin{thebibliography}{9}

\bibitem{trust-systems}
Trust systems: the certificate lemma and transport of trust.
\newblock Companion formal development.

\end{thebibliography}

\end{document}
